% GENERATED FILE. Edit canonical lessons, then run npm run generate:books.
% canonical-source-hash: fnv1a64:9a75588a67ec2b41
% canonical-lessons: ML-C375-raksakan, ML-C375-sannaddhapravarttakan, ML-C375-raksikkuka, ML-C375-raksappetuka, ML-C375-kutunnuka

\chapter{Rescue}
\label{ch:ml-rescue}
\hlchaptermodality{375}

\begin{chapteropening}
\textbf{By the end of this chapter:} \emph{I can say a rescuer, a volunteer, to rescue, to escape and to get stuck.}

Complete the last lesson of chapter 375: I can say a rescuer, a volunteer, to rescue, to escape and to get stuck.
\end{chapteropening}

\section[\texorpdfstring{\ml{രക്ഷകൻ} (rakṣakan)}{രക്ഷകൻ (rakṣakan)}]{\ml{രക്ഷകൻ} (rakṣakan) — a rescuer; a saviour}

\label{lesson:ML-C375-raksakan}



\begin{quote}
Before the new one: say the Malayalam for a state of emergency, then the Malayalam for to evacuate.
\end{quote}

\subsection*{You'll want to know: \ml{രക്ഷകൻ}}
\textbf{\ml{രക്ഷകൻ}} — \emph{rakṣakan} — \textquotedblleft{}a rescuer; a saviour\textquotedblright{}.

\begin{quote}
\ml{ആ} \ml{ചെറുപ്പക്കാരൻ} \ml{ഗ്രാമത്തിന്റെ} \ml{രക്ഷകനായി}.
\end{quote}

\begin{quote}
\emph{ā ceṟuppakkāran grāmattinṟe rakṣakanāyi.}
\end{quote}

\begin{quote}
That young man became the village's saviour.
\end{quote}

The noun \ml{രക്ഷ} means protection; a woman rescuer is \ml{രക്ഷക}.

\begin{cousinweb}
\textbf{\ml{രക്ഷകൻ}} (\emph{rakṣakan}) is borrowed from Sanskrit \emph{rakṣaka}, \textquotedblleft{}protecting, guarding\textquotedblright{} (Monier-Williams), with the Malayalam masculine ending -\ml{ൻ}.
\end{cousinweb}

\begin{grammarlens}[title={What you've built}]
One more word for telling what happened.
\end{grammarlens}

\subsection*{Your turn}
\begin{itemize}
\raggedright
  \item \emph{Say it:} \emph{rakṣakan}
  \item \emph{Say it:} \emph{rakṣakan}, once more
  \item \emph{Say it:} \emph{ā ceṟuppakkāran grāmattinṟe rakṣakanāyi.} — That young man became the village's saviour.
  \item \emph{Recall:} say the Malayalam for a state of emergency, then the Malayalam for to evacuate, then say \emph{rakṣakan} again
\end{itemize}

\begin{tcolorbox}[breakable,colback=teal!4,colframe=teal!35!black,title={Before you move on}]
What does \ml{രക്ഷകൻ} mean? (\textquotedblleft{}A rescuer; a saviour\textquotedblright{}.) Say it once more.
\end{tcolorbox}
\section[\texorpdfstring{\ml{സന്നദ്ധപ്രവർത്തകൻ} (sannaddhapravarttakan)}{സന്നദ്ധപ്രവർത്തകൻ (sannaddhapravarttakan)}]{\ml{സന്നദ്ധപ്രവർത്തകൻ} (sannaddhapravarttakan) — a volunteer}

\label{lesson:ML-C375-sannaddhapravarttakan}



\begin{quote}
Before the new one: say the Malayalam for to evacuate, then the Malayalam for a rescuer.
\end{quote}

\subsection*{You'll want to know: \ml{സന്നദ്ധപ്രവർത്തകൻ}}
\textbf{\ml{സന്നദ്ധപ്രവർത്തകൻ}} — \emph{sannaddhapravarttakan} — \textquotedblleft{}a volunteer\textquotedblright{}.

\begin{quote}
\ml{സന്നദ്ധപ്രവർത്തകർ} \ml{ആളുകൾക്ക്} \ml{ഭക്ഷണം} \ml{കൊടുത്തു}.
\end{quote}

\begin{quote}
\emph{sannaddhapravarttakar āḷukaḷkkŭ bhakṣaṇaṁ koṭuttu.}
\end{quote}

\begin{quote}
Volunteers gave people food.
\end{quote}

Plural: \ml{സന്നദ്ധപ്രവർത്തകർ}, \emph{sannaddhapravarttakar}.

\ml{സന്നദ്ധ}, ready, + \ml{പ്രവർത്തകൻ}, a worker: one ready to work without pay.

\begin{cousinweb}
\textbf{\ml{സന്നദ്ധപ്രവർത്തകൻ}} (\emph{sannaddhapravarttakan}) is built from Sanskrit \emph{saṁnaddha}, \textquotedblleft{}girt, bound, armed\textquotedblright{} (CDIAL 12923), so ready for action, and \ml{പ്രവർത്തകൻ}, \textquotedblleft{}a worker\textquotedblright{}, from \emph{pravartaka} (Monier-Williams).
\end{cousinweb}

\begin{grammarlens}[title={What you've built}]
One more word for telling what happened.
\end{grammarlens}

\subsection*{Your turn}
\begin{itemize}
\raggedright
  \item \emph{Say it:} \emph{sannaddhapravarttakan}
  \item \emph{Say it:} \emph{sannaddhapravarttakan}, once more
  \item \emph{Say it:} \emph{sannaddhapravarttakar āḷukaḷkkŭ bhakṣaṇaṁ koṭuttu.} — Volunteers gave people food.
  \item \emph{Recall:} say the Malayalam for to evacuate, then the Malayalam for a rescuer, then say \emph{sannaddhapravarttakan} again
\end{itemize}

\begin{tcolorbox}[breakable,colback=teal!4,colframe=teal!35!black,title={Before you move on}]
What does \ml{സന്നദ്ധപ്രവർത്തകൻ} mean? (\textquotedblleft{}A volunteer\textquotedblright{}.) Say it once more.
\end{tcolorbox}
\section[\texorpdfstring{\ml{രക്ഷിക്കുക} (rakṣikkuka)}{രക്ഷിക്കുക (rakṣikkuka)}]{\ml{രക്ഷിക്കുക} (rakṣikkuka) — to rescue; to save}

\label{lesson:ML-C375-raksikkuka}



\begin{quote}
Before the new one: say the Malayalam for a rescuer, then the Malayalam for a volunteer.
\end{quote}

\subsection*{You'll want to know: \ml{രക്ഷിക്കുക}}
\textbf{\ml{രക്ഷിക്കുക}} — \emph{rakṣikkuka} — \textquotedblleft{}to rescue; to save\textquotedblright{}.

\begin{quote}
\ml{നാട്ടുകാർ} \ml{പുഴയിൽ} \ml{വീണ} \ml{കുട്ടിയെ} \ml{രക്ഷിച്ചു}.
\end{quote}

\begin{quote}
\emph{nāṭṭukār puḻayil vīṇa kuṭṭiye rakṣiccu.}
\end{quote}

\begin{quote}
The local people rescued the child who fell into the river.
\end{quote}

Past: \ml{രക്ഷിച്ചു}, \emph{rakṣiccu}.

The one saved takes -\ml{എ}. Help! is \ml{രക്ഷിക്കണേ}!

\begin{cousinweb}
\textbf{\ml{രക്ഷിക്കുക}} (\emph{rakṣikkuka}) is built from Sanskrit \emph{rakṣā}, \textquotedblleft{}protection, preservation\textquotedblright{} (CDIAL 10551), with the Malayalam verb ending -\ml{ിക്കുക}.
\end{cousinweb}

\begin{grammarlens}[title={What you've built}]
One more word for telling what happened.
\end{grammarlens}

\subsection*{Your turn}
\begin{itemize}
\raggedright
  \item \emph{Say it:} \emph{rakṣikkuka}
  \item \emph{Say it:} \emph{rakṣikkuka}, once more
  \item \emph{Say it:} \emph{nāṭṭukār puḻayil vīṇa kuṭṭiye rakṣiccu.} — The local people rescued the child who fell into the river.
  \item \emph{Recall:} say the Malayalam for a rescuer, then the Malayalam for a volunteer, then say \emph{rakṣikkuka} again
\end{itemize}

\begin{tcolorbox}[breakable,colback=teal!4,colframe=teal!35!black,title={Before you move on}]
What does \ml{രക്ഷിക്കുക} mean? (\textquotedblleft{}To rescue; to save\textquotedblright{}.) Say it once more.
\end{tcolorbox}
\section[\texorpdfstring{\ml{രക്ഷപ്പെടുക} (rakṣappeṭuka)}{രക്ഷപ്പെടുക (rakṣappeṭuka)}]{\ml{രക്ഷപ്പെടുക} (rakṣappeṭuka) — to escape; to be saved}

\label{lesson:ML-C375-raksappetuka}



\begin{quote}
Before the new one: say the Malayalam for a volunteer, then the Malayalam for to rescue.
\end{quote}

\subsection*{You'll want to know: \ml{രക്ഷപ്പെടുക}}
\textbf{\ml{രക്ഷപ്പെടുക}} — \emph{rakṣappeṭuka} — \textquotedblleft{}to escape; to be saved\textquotedblright{}.

\begin{quote}
\ml{തീപിടിത്തത്തിൽ} \ml{നിന്ന്} \ml{എല്ലാവരും} \ml{രക്ഷപ്പെട്ടു}.
\end{quote}

\begin{quote}
\emph{tīpiṭittattil ninnŭ ellāvaruṁ rakṣappeṭṭu.}
\end{quote}

\begin{quote}
Everyone escaped from the fire.
\end{quote}

Past: \ml{രക്ഷപ്പെട്ടു}, \emph{rakṣappeṭṭu}.

From what takes -\ml{ഇൽ} \ml{നിന്ന്}. -\ml{പ്പെടുക} makes the saving happen to the subject.

\begin{cousinweb}
\textbf{\ml{രക്ഷപ്പെടുക}} (\emph{rakṣappeṭuka}) is built from Sanskrit \emph{rakṣā}, \textquotedblleft{}protection\textquotedblright{} (CDIAL 10551), with -\ml{പ്പെടുക}, \textquotedblleft{}to come to be\textquotedblright{}: to come into safety.
\end{cousinweb}

\begin{grammarlens}[title={What you've built}]
One more word for telling what happened.
\end{grammarlens}

\subsection*{Your turn}
\begin{itemize}
\raggedright
  \item \emph{Say it:} \emph{rakṣappeṭuka}
  \item \emph{Say it:} \emph{rakṣappeṭuka}, once more
  \item \emph{Say it:} \emph{tīpiṭittattil ninnŭ ellāvaruṁ rakṣappeṭṭu.} — Everyone escaped from the fire.
  \item \emph{Recall:} say the Malayalam for a volunteer, then the Malayalam for to rescue, then say \emph{rakṣappeṭuka} again
\end{itemize}

\begin{tcolorbox}[breakable,colback=teal!4,colframe=teal!35!black,title={Before you move on}]
What does \ml{രക്ഷപ്പെടുക} mean? (\textquotedblleft{}To escape; to be saved\textquotedblright{}.) Say it once more.
\end{tcolorbox}
\section[\texorpdfstring{\ml{കുടുങ്ങുക} (kuṭuṅṅuka)}{കുടുങ്ങുക (kuṭuṅṅuka)}]{\ml{കുടുങ്ങുക} (kuṭuṅṅuka) — to get stuck; to be trapped}

\label{lesson:ML-C375-kutunnuka}



\begin{quote}
Before the new one: say the Malayalam for to rescue, then the Malayalam for to escape.
\end{quote}

\subsection*{You'll want to know: \ml{കുടുങ്ങുക}}
\textbf{\ml{കുടുങ്ങുക}} — \emph{kuṭuṅṅuka} — \textquotedblleft{}to get stuck; to be trapped\textquotedblright{}.

\begin{quote}
\ml{മൂന്ന്} \ml{പേർ} \ml{കെട്ടിടത്തിനുള്ളിൽ} \ml{കുടുങ്ങി}.
\end{quote}

\begin{quote}
\emph{mūnnŭ pēr keṭṭiṭattinuḷḷil kuṭuṅṅi.}
\end{quote}

\begin{quote}
Three people were trapped inside the building.
\end{quote}

Past: \ml{കുടുങ്ങി}, \emph{kuṭuṅṅi}.

Also of traffic: \ml{ട്രാഫിക്കിൽ} \ml{കുടുങ്ങി}, stuck in traffic.

\begin{cousinweb}
\textbf{\ml{കുടുങ്ങുക}} (\emph{kuṭuṅṅuka}) is a native Malayalam verb whose origin is not settled. It is said of people, vehicles and things caught where they cannot move.
\end{cousinweb}

\begin{grammarlens}[title={What you've built}]
One more word for telling what happened.
\end{grammarlens}

\subsection*{Your turn}
\begin{itemize}
\raggedright
  \item \emph{Say it:} \emph{kuṭuṅṅuka}
  \item \emph{Say it:} \emph{kuṭuṅṅuka}, once more
  \item \emph{Say it:} \emph{mūnnŭ pēr keṭṭiṭattinuḷḷil kuṭuṅṅi.} — Three people were trapped inside the building.
  \item \emph{Recall:} say the Malayalam for to rescue, then the Malayalam for to escape, then say \emph{kuṭuṅṅuka} again
\end{itemize}

\begin{tcolorbox}[breakable,colback=teal!4,colframe=teal!35!black,title={Before you move on}]
What does \ml{കുടുങ്ങുക} mean? (\textquotedblleft{}To get stuck; to be trapped\textquotedblright{}.) Say it once more.
\end{tcolorbox}
